\CWHeader{I. Персистентные структуры данных}

\CWProblem{
Дан набор горизонтальных отрезков, и набор точек. Для каждой точки определить сколько отрезков лежит строго над ней.

Решение должно работать online, то есть должно обрабатывать запросы по одному после построения необходимой структуры
данных по входным данным. Чтение входных данных и запросов вместе и построение по ним общей структуры запрещено.

\textbf{Формат ввода}

В первой строке вам даны два числа $n$ и $m$ ($1 \le n, m \le 10^5$) — количество отрезков и количество точек соответственно.
В следующих $n$ строках вам заданы отрезки, в виде троек чисел $l$, $r$ и $h$ ($-10^9 \le l < r \le 10^9, -10^9 \le g \le 10^9$)
— координаты $x$ левой и правой границ отрезка и координата $y$ отрезка соответственно. В следующих $m$ строках вам даны
пары чисел $x$, $y$ ($-10^9 \le x, y \le 10^9$) — координаты точек.

\textbf{Формат вывода}

Для каждой точки выведите количество отрезков над ней.

}
\pagebreak
